\documentclass[12pt,a4paper]{article}
\usepackage[utf8]{inputenc}
\usepackage[T5]{fontenc}
\usepackage[vietnamese]{babel}
\usepackage{amsmath, amssymb}
\usepackage{array}
\usepackage{enumitem}
\usepackage{setspace}
\setstretch{1.18}
\usepackage{graphicx}
\usepackage{ifthen}
\usepackage{tikz}
\usetikzlibrary{arrows.meta, positioning, patterns}
\IfFileExists{pgf-pie.sty}{\usepackage{pgf-pie}}{}
\usepackage{pgfplots}
\pgfplotsset{compat=1.18}
\usepackage{chemfig}
\usepackage{longtable,booktabs}
\usepackage{adjustbox}
\providecommand{\errorpart}[2]{$\underset{\textbf{#1}}{\underline{\text{#2}}}$}
\usepackage[top=0.8in, left=0.75in, right=0.75in, bottom=0.85in, headheight=38pt, headsep=18pt, footskip=1.2cm]{geometry}
\usepackage{fancyhdr}
\usepackage[colorlinks=true, linkcolor=blue, urlcolor=blue]{hyperref}

\setlength{\parindent}{0pt}
\setlength{\parskip}{5pt plus 2pt minus 1pt}

\setlist[enumerate]{itemsep=4pt, parsep=2pt, topsep=4pt, partopsep=2pt}
\setlist[itemize]{itemsep=4pt, parsep=2pt, topsep=4pt}

\makeatletter
\AtEndDocument{\immediate\write\@auxout{\string\newlabel{LecturingLastPage}{{}{\thepage}{}{}{}}}}
\AtEndDocument{\immediate\write\@auxout{\string\newlabel{LecturingPageCount}{{}{\two@digits{\value{page}}}{}{}{}}}}
\makeatother
\pagestyle{fancy}
\fancyhf{}
\cfoot{\small Trang \thepage/\pageref*{LecturingLastPage} - Mã đề thi 2615}
\renewcommand{\headrulewidth}{0pt}
\renewcommand{\footrulewidth}{0pt}

\begin{document}

\sloppy

\noindent
\begin{minipage}[t]{0.37\textwidth}
    \small
    \centering
    \IfFileExists{logo.png}{\includegraphics[height=0.9cm]{logo.png}\\[0.1cm]}{%
      \IfFileExists{logo.jpg}{\includegraphics[height=0.9cm]{logo.jpg}\\[0.1cm]}{}%
    }
    {\bfseries \mbox{}}\\[-0.2cm]
    \rule{3cm}{0.5pt}\\[0.1cm]
    {\bfseries ĐỀ LUYỆN TẬP}\\{\small (Đề thi có \pageref*{LecturingPageCount} trang)}
\end{minipage}%
\hfill
\begin{minipage}[t]{0.62\textwidth}
    \small
    \centering
    {\bfseries \sbox0{KỲ THI TỐT NGHIỆP TRUNG HỌC PHỔ THÔNG NĂM 2026}\ifdim\wd0>\linewidth KỲ THI TỐT NGHIỆP TRUNG HỌC PHỔ~THÔNG\\NĂM 2026\else\usebox0\fi}\\[0.15cm]
    {\bfseries Môn thi: HOÁ HỌC}\\[0.15cm]
    {\itshape Thời gian làm bài: 50 phút, không kể thời gian phát đề}
\end{minipage}

\vspace{0.4cm}
\noindent
\begin{minipage}[t]{0.68\textwidth}
    Họ, tên thí sinh: \dotfill\\[0.15cm]
    Số báo danh: \dotfill
\end{minipage}%
\hfill
\fbox{\begin{minipage}{3.6cm}\centering\textbf{Mã đề: 2615}\end{minipage}}

\vspace{0.5cm}

% LECTURING-PART: P1
\textbf{Phần I: Trắc nghiệm nhiều lựa chọn} \textit{Thí sinh trả lời từ câu 1 đến câu 18. Mỗi câu hỏi thí sinh chỉ chọn một phương án.}

\noindent \textbf{Câu 1:} Kim loại nào sau đây tác dụng mạnh với nước ở nhiệt độ thường, tạo thành dung dịch base và khí hydrogen?\\
\noindent\makebox[0.24\textwidth][l]{A. $\text{Ag}$.}\makebox[0.24\textwidth][l]{B. $\text{Cu}$.}\makebox[0.24\textwidth][l]{C. $\text{Na}$.}\makebox[0.24\textwidth][l]{D. $\text{Fe}$.}

\noindent \textbf{Câu 2:} Chất nào sau đây là amine bậc hai?\\
\noindent A. $\text{CH}_3\text{NH}_2$. \hfill B. $(\text{CH}_3)_2\text{NH}$. \hfill C. $(\text{CH}_3)_3\text{N}$. \hfill D. $\text{C}_6\text{H}_5\text{NH}_2$.\\

\noindent \textbf{Câu 3:} Công thức cấu tạo dạng mạch vòng của monosaccharide X được biểu diễn ở hình bên dưới.\\
\begin{center}
\begin{adjustbox}{max width=0.95\linewidth, max totalheight=0.38\textheight}
\setchemfig{atom sep=2.4em}\chemfig{O*5(-(-[::25]OH)(-[::-45]CH_2OH)-(-OH)-(-OH)-(-CH_2OH)-)}
\end{adjustbox}
\end{center}
Chất X là\\
\noindent\makebox[0.24\textwidth][l]{A. fructose.}\makebox[0.24\textwidth][l]{B. maltose.}\makebox[0.24\textwidth][l]{C. saccharose.}\makebox[0.24\textwidth][l]{D. glucose.}

\noindent \textbf{Câu 4:} Cho bảng giá trị thế điện cực chuẩn của một số cặp oxi hoá -- khử ở $25\,^{\circ}\text{C}$ như sau:\\
\begin{center}
\begin{adjustbox}{max width=\linewidth}
\begin{tabular}{|c|c|c|c|c|}
\hline
Cặp oxi hoá -- khử & $\text{Mg}^{2+}/\text{Mg}$ & $\text{Zn}^{2+}/\text{Zn}$ & $\text{Cu}^{2+}/\text{Cu}$ & $\text{Ag}^{+}/\text{Ag}$ \\
\hline
$E^{\circ}$ (V) & $-2{,}37$ & $-0{,}76$ & $+0{,}34$ & $+0{,}80$ \\
\hline
\end{tabular}
\end{adjustbox}
\end{center}
Sức điện động chuẩn của pin điện hoá $\text{Zn}$ -- $\text{Ag}$ là\\
\noindent\makebox[0.24\textwidth][l]{A. $0{,}04$ V.}\makebox[0.24\textwidth][l]{B. $1{,}10$ V.}\makebox[0.24\textwidth][l]{C. $1{,}56$ V.}\makebox[0.24\textwidth][l]{D. $2{,}36$ V.}

\noindent \textbf{Câu 5:} Thuỷ phân hoàn toàn ethyl acetate trong dung dịch $\text{NaOH}$ đun nóng, thu được các sản phẩm hữu cơ là\\
\noindent A. $\text{CH}_3\text{COONa}$ và $\text{C}_2\text{H}_5\text{OH}$. \hfill B. $\text{C}_2\text{H}_5\text{COONa}$ và $\text{CH}_3\text{OH}$.\\
\noindent C. $\text{HCOONa}$ và $\text{C}_3\text{H}_7\text{OH}$. \hfill D. $\text{CH}_3\text{COOH}$ và $\text{C}_2\text{H}_5\text{ONa}$.\\

\noindent \textbf{Câu 6:} Hydrogen hoá hoàn toàn triolein (xúc tác $\text{Ni}$, đun nóng, áp suất cao) thu được chất béo nào sau đây?\\
\noindent\makebox[0.24\textwidth][l]{A. Trilinolein.}\makebox[0.24\textwidth][l]{B. Tristearin.}\makebox[0.24\textwidth][l]{C. Tripalmitin.}\makebox[0.24\textwidth][l]{D. Glycerol.}

\noindent \textbf{Câu 7:} Polymer nào sau đây được điều chế bằng phản ứng trùng ngưng?\\
\noindent\makebox[0.48\textwidth][l]{A. Polystyrene.}\makebox[0.48\textwidth][l]{B. Poly(vinyl chloride).}\\
\noindent\makebox[0.48\textwidth][l]{C. Nylon-6,6.}\makebox[0.48\textwidth][l]{D. Polyethylene.}

\noindent \textbf{Câu 8:} Dung dịch của chất nào sau đây làm quỳ tím chuyển sang màu xanh?\\
\noindent\makebox[0.24\textwidth][l]{A. Lysine.}\makebox[0.24\textwidth][l]{B. Glutamic acid.}\makebox[0.24\textwidth][l]{C. Glycine.}\makebox[0.24\textwidth][l]{D. Alanine.}

\noindent \textbf{Câu 9:} Nhỏ từ từ đến dư dung dịch $\text{NH}_3$ vào ống nghiệm đựng dung dịch $\text{CuSO}_4$. Hiện tượng quan sát được là\\
\noindent A. xuất hiện kết tủa màu xanh, sau đó kết tủa tan tạo dung dịch màu xanh lam đậm.\\
\noindent B. xuất hiện kết tủa màu xanh, kết tủa không tan khi thêm $\text{NH}_3$ dư.\\
\noindent C. xuất hiện kết tủa màu trắng, sau đó kết tủa tan tạo dung dịch không màu.\\
\noindent D. có khí mùi khai thoát ra và dung dịch chuyển sang không màu.

\noindent \textbf{Câu 10:} Nguyên tử của nguyên tố X có số hiệu nguyên tử $Z = 17$. Vị trí của X trong bảng tuần hoàn là\\
\noindent\makebox[0.48\textwidth][l]{A. chu kì 3, nhóm VA.}\makebox[0.48\textwidth][l]{B. chu kì 4, nhóm IA.}\\
\noindent\makebox[0.48\textwidth][l]{C. chu kì 2, nhóm VIIA.}\makebox[0.48\textwidth][l]{D. chu kì 3, nhóm VIIA.}

\noindent \textbf{Câu 11:} Điện phân dung dịch $\text{CuSO}_4$ với các điện cực trơ. Quá trình xảy ra ở anode là\\
\noindent A. sự oxi hoá phân tử $\text{H}_2\text{O}$.\\
\noindent B. sự oxi hoá ion $\text{SO}_4^{2-}$.\\
\noindent C. sự khử ion $\text{Cu}^{2+}$.\\
\noindent D. sự khử phân tử $\text{H}_2\text{O}$.

\noindent \textbf{Câu 12:} Một mẫu nước giếng chứa các ion $\text{Ca}^{2+}$, $\text{Mg}^{2+}$, $\text{Cl}^{-}$ và $\text{SO}_4^{2-}$. Hoá chất nào sau đây làm mềm được mẫu nước này?\\
\noindent\makebox[0.48\textwidth][l]{A. $\text{Ca(OH)}_2$.}\makebox[0.48\textwidth][l]{B. $\text{Na}_2\text{CO}_3$.}\\
\noindent\makebox[0.48\textwidth][l]{C. $\text{HCl}$.}\makebox[0.48\textwidth][l]{D. $\text{NaCl}$.}

\noindent \textbf{Câu 13:} Một ống thép (thành phần chính là $\text{Fe}$) được nối với một thanh kim loại M rồi cùng đặt trong đất ẩm. Với kim loại M nào sau đây thì ống thép bị ăn mòn điện hoá nhanh hơn?\\
\begin{center}
\begin{adjustbox}{max width=0.8\linewidth, max totalheight=0.25\textheight}
\begin{tikzpicture}[scale=0.9]
  \fill[brown!25] (0,0) rectangle (9,3);
  \draw[thick] (0,0) rectangle (9,3);
  \node[fill=brown!25, inner sep=2pt] at (4.5,0.45) {đất ẩm};
  \fill[gray!55] (0.8,1.2) rectangle (4.2,1.9);
  \draw (0.8,1.2) rectangle (4.2,1.9);
  \node[above] at (2.5,1.95) {ống thép (Fe)};
  \fill[gray!30] (6.4,1.1) rectangle (7.6,2.0);
  \draw (6.4,1.1) rectangle (7.6,2.0);
  \node[above] at (7.0,2.05) {kim loại M};
  \draw[thick] (4.2,1.55) -- (6.4,1.55);
  \node[below, font=\small] at (5.3,1.5) {dây dẫn};
\end{tikzpicture}
\end{adjustbox}
\end{center}
\noindent\makebox[0.24\textwidth][l]{A. $\text{Cu}$.}\makebox[0.24\textwidth][l]{B. $\text{Mg}$.}\makebox[0.24\textwidth][l]{C. $\text{Zn}$.}\makebox[0.24\textwidth][l]{D. $\text{Al}$.}

\noindent \textbf{Câu 14:} Xà phòng hoá hoàn toàn $8{,}8$ gam ethyl acetate bằng dung dịch $\text{NaOH}$ vừa đủ, cô cạn dung dịch sau phản ứng thu được $m$ gam muối khan. Cho nguyên tử khối: $\text{H} = 1$; $\text{C} = 12$; $\text{O} = 16$; $\text{Na} = 23$. Giá trị của $m$ là\\
\noindent\makebox[0.24\textwidth][l]{A. $4{,}6$.}\makebox[0.24\textwidth][l]{B. $6{,}8$.}\makebox[0.24\textwidth][l]{C. $8{,}2$.}\makebox[0.24\textwidth][l]{D. $9{,}6$.}

\noindent \textbf{Câu 15:} Lên men $36$ gam glucose thành ethanol với hiệu suất $80\%$, toàn bộ khí $\text{CO}_2$ sinh ra được thu lại. Biết ở điều kiện chuẩn ($25\,^{\circ}\text{C}$, $1$ bar) một mol khí chiếm thể tích $24{,}79$ lít; nguyên tử khối: $\text{H} = 1$; $\text{C} = 12$; $\text{O} = 16$. Thể tích khí $\text{CO}_2$ thu được ở điều kiện chuẩn là\\
\noindent\makebox[0.24\textwidth][l]{A. $3{,}97$ lít.}\makebox[0.24\textwidth][l]{B. $4{,}96$ lít.}\makebox[0.24\textwidth][l]{C. $7{,}93$ lít.}\makebox[0.24\textwidth][l]{D. $9{,}92$ lít.}

\noindent \textbf{Câu 16:} Cho $4{,}5$ gam ethylamine ($\text{C}_2\text{H}_5\text{NH}_2$) tác dụng vừa đủ với dung dịch $\text{HCl}$, thu được dung dịch chứa $m$ gam muối. Cho nguyên tử khối: $\text{H} = 1$; $\text{C} = 12$; $\text{N} = 14$; $\text{Cl} = 35{,}5$. Giá trị của $m$ là\\
\noindent A. $7{,}65$. \hfill B. $8{,}15$. \hfill C. $8{,}55$. \hfill D. $9{,}15$.\\

\noindent \textbf{Câu 17:} Phát biểu nào sau đây về protein là \textbf{đúng}?\\
\noindent A. Protein đơn giản khi thuỷ phân hoàn toàn chỉ thu được các $\alpha$-amino acid.\\
\noindent B. Protein không bị đông tụ khi đun nóng.\\
\noindent C. Liên kết peptide là liên kết giữa nhóm $\text{-COOH}$ và nhóm $\text{-OH}$.\\
\noindent D. Tất cả các protein đều tan tốt trong nước tạo dung dịch keo.

\noindent \textbf{Câu 18:} Hoà tan hoàn toàn $7{,}8$ gam hỗn hợp gồm $\text{Mg}$ và $\text{Al}$ trong dung dịch $\text{HCl}$ dư, thu được $9{,}916$ lít khí $\text{H}_2$ ở điều kiện chuẩn (một mol khí chiếm $24{,}79$ lít). Cho nguyên tử khối: $\text{Mg} = 24$; $\text{Al} = 27$. Khối lượng $\text{Al}$ trong hỗn hợp ban đầu là\\
\noindent\makebox[0.24\textwidth][l]{A. $2{,}4$ gam.}\makebox[0.24\textwidth][l]{B. $2{,}7$ gam.}\makebox[0.24\textwidth][l]{C. $4{,}05$ gam.}\makebox[0.24\textwidth][l]{D. $5{,}4$ gam.}

\vspace{0.5cm}
% LECTURING-PART: P2
\textbf{Phần II: Trắc nghiệm đúng/sai} \textit{Thí sinh trả lời từ câu 1 đến câu 4. Trong mỗi ý a), b), c), d) ở mỗi câu, thí sinh chọn đúng hoặc sai.}

\noindent \textbf{Câu 1:} Một nhóm học sinh tiến hành điện phân dung dịch $\text{CuSO}_4$ với hai điện cực graphite theo sơ đồ ở hình bên dưới, cường độ dòng điện không đổi $I = 1{,}93$ A trong thời gian $2000$ giây. Biết hiệu suất điện phân là $100\%$, hằng số Faraday $F = 96485$ C/mol, nguyên tử khối $\text{Cu} = 64$, và trong suốt thời gian điện phân ion $\text{Cu}^{2+}$ vẫn còn dư.\\
\begin{center}
\begin{adjustbox}{max width=0.95\linewidth, max totalheight=0.21\textheight}
\begin{tikzpicture}[scale=0.9]
  % cốc
  \draw[thick] (0,4) -- (0,0) -- (6,0) -- (6,4);
  % mặt dung dịch
  \fill[gray!15] (0.02,0.02) rectangle (5.98,3);
  \draw (0,3) -- (6,3);
  % điện cực
  \fill[gray!70] (1.3,1.0) rectangle (1.7,4.6);
  \draw (1.3,1.0) rectangle (1.7,4.6);
  \fill[gray!70] (4.3,1.0) rectangle (4.7,4.6);
  \draw (4.3,1.0) rectangle (4.7,4.6);
  % dây dẫn và nguồn
  \draw (1.5,4.6) -- (1.5,5.8) -- (2.85,5.8);
  \draw (4.5,4.6) -- (4.5,5.8) -- (3.15,5.8);
  \draw[very thick] (2.85,5.55) -- (2.85,6.05); % vạch ngắn: cực âm
  \draw[very thick] (3.15,5.35) -- (3.15,6.25); % vạch dài: cực dương


  % nhãn

  \draw (-0.1,3.8) -- (1.3,3.8);

  \draw (6.1,3.8) -- (4.7,3.8);
    % captions last so nothing paints over them
    % captions last so nothing paints over them
    \node[above] at (2.6,6.0) {$-$};
\node[above] at (3.4,6.2) {$+$};
\node[left] at (-0.1,3.8) {cathode};
\node[right] at (6.1,3.8) {anode};
\node[fill=gray!15, inner sep=2pt] at (3,0.5) {dung dịch $\text{CuSO}_4$};
\end{tikzpicture}
\end{adjustbox}
\end{center}
\noindent a) Trong quá trình điện phân, ở điện cực nối với cực âm của nguồn có lớp kim loại màu đỏ bám vào.\\

\noindent b) Khí thoát ra ở anode là khí hydrogen.\\

\noindent c) Khối lượng kim loại bám vào cathode sau $2000$ giây xấp xỉ $1{,}28$ gam.\\

\noindent d) Giá trị pH của dung dịch tăng dần trong quá trình điện phân.\\

\noindent \textbf{Câu 2:} Tiến hành thí nghiệm phản ứng tráng bạc của glucose theo các bước sau:\\
Bước 1: Cho $2$ mL dung dịch $\text{AgNO}_3$ $1\%$ vào ống nghiệm sạch, nhỏ từ từ dung dịch $\text{NH}_3$ $5\%$ vào, lắc đều đến khi kết tủa vừa tạo thành tan hết.\\
Bước 2: Thêm tiếp $1$ mL dung dịch glucose $10\%$ vào ống nghiệm.\\
Bước 3: Đặt ống nghiệm vào cốc nước nóng khoảng $60$ -- $70\,^{\circ}\text{C}$ trong vài phút.\\
Cho nguyên tử khối: $\text{H} = 1$; $\text{C} = 12$; $\text{O} = 16$; $\text{Ag} = 108$.\\
\noindent a) Sau bước 3, trên thành ống nghiệm xuất hiện một lớp kim loại sáng như gương.\\

\noindent b) Trong phản ứng ở bước 3, glucose đóng vai trò chất oxi hoá.\\

\noindent c) Nếu thay dung dịch glucose ở bước 2 bằng dung dịch saccharose thì sau bước 3 cũng thu được lớp kim loại sáng như gương.\\

\noindent d) Nếu lượng glucose phản ứng là $0{,}9$ gam và phản ứng xảy ra hoàn toàn thì khối lượng kim loại tạo thành là $1{,}08$ gam.\\

\noindent \textbf{Câu 3:} Tiến hành thí nghiệm xà phòng hoá chất béo theo các bước sau:\\
Bước 1: Cho khoảng $2$ gam mỡ lợn và $5$ mL dung dịch $\text{NaOH}$ $40\%$ vào bát sứ.\\
Bước 2: Đun sôi nhẹ hỗn hợp, liên tục khuấy đều bằng đũa thuỷ tinh, thỉnh thoảng thêm vài giọt nước cất.\\
Bước 3: Sau khoảng $8$ -- $10$ phút, rót thêm vào hỗn hợp $10$ mL dung dịch $\text{NaCl}$ bão hoà nóng, khuấy nhẹ rồi để nguội.\\
Cho nguyên tử khối: $\text{H} = 1$; $\text{C} = 12$; $\text{O} = 16$; $\text{Na} = 23$.\\
\noindent a) Sau bước 3, có một lớp chất rắn màu trắng nổi lên trên bề mặt hỗn hợp.\\

\noindent b) Mục đích chính của việc thêm dung dịch $\text{NaCl}$ bão hoà ở bước 3 là làm tăng tốc độ phản ứng xà phòng hoá.\\

\noindent c) Việc thêm nước cất ở bước 2 giúp hỗn hợp không bị cạn khô, nhờ đó phản ứng xà phòng hoá tiếp tục xảy ra.\\

\noindent d) Xà phòng hoá hoàn toàn $8{,}9$ gam tristearin bằng dung dịch $\text{NaOH}$ dư thu được $2{,}76$ gam glycerol.\\

\noindent \textbf{Câu 4:} Để xác định nồng độ của một dung dịch $\text{NaOH}$, một học sinh lấy chính xác $10{,}0$ mL dung dịch này cho vào bình tam giác, thêm 2 giọt phenolphthalein, rồi chuẩn độ bằng dung dịch $\text{HCl}$ $0{,}100$ M. Kết quả ba lần chuẩn độ được ghi ở bảng sau:\\
\begin{center}
\begin{adjustbox}{max width=\linewidth}
\begin{tabular}{|c|c|c|c|}
\hline
Lần chuẩn độ & 1 & 2 & 3 \\
\hline
Thể tích dung dịch $\text{HCl}$ đã dùng (mL) & $10{,}1$ & $9{,}9$ & $10{,}0$ \\
\hline
\end{tabular}
\end{adjustbox}
\end{center}
\noindent a) Thời điểm kết thúc chuẩn độ là khi dung dịch trong bình tam giác chuyển từ màu hồng sang không màu và giữ nguyên trong khoảng 20 giây.\\

\noindent b) Thể tích trung bình của dung dịch $\text{HCl}$ đã dùng trong ba lần chuẩn độ là $10{,}0$ mL.\\

\noindent c) Nồng độ của dung dịch $\text{NaOH}$ ban đầu là $0{,}010$ M.\\

\noindent d) Dung dịch $\text{HCl}$ trong thí nghiệm này được chứa trong pipette và nhỏ từ từ vào bình tam giác.\\

\vspace{0.5cm}
% LECTURING-PART: P3
\textbf{Phần III: Trắc nghiệm trả lời ngắn} \textit{Thí sinh trả lời từ câu 1 đến câu 6. Thí sinh chỉ điền một con số (có thể có dấu âm, dấu phẩy thập phân) vào ô trống.}

\noindent \textbf{Câu 1:} Cho các polymer sau: (1) polyethylene; (2) nylon-6,6; (3) cellulose; (4) poly(vinyl chloride); (5) polyacrylonitrile. Liệt kê các polymer được điều chế bằng phản ứng trùng hợp theo thứ tự số thứ tự tăng dần (ví dụ: 123, 245).\\

\noindent \textbf{Câu 2:} Có bao nhiêu ester mạch hở có cùng công thức phân tử $\text{C}_4\text{H}_8\text{O}_2$? \textit{(Làm tròn đến chữ số hàng đơn vị)}\\

\noindent \textbf{Câu 3:} Trong công nghiệp, nhôm được sản xuất bằng cách điện phân nóng chảy $\text{Al}_2\text{O}_3$. Từ $10$ tấn quặng bauxite chứa $51\%$ khối lượng $\text{Al}_2\text{O}_3$ (phần còn lại không chứa nhôm), với hiệu suất của cả quá trình là $90\%$, thu được bao nhiêu tấn nhôm? Cho nguyên tử khối: $\text{O} = 16$; $\text{Al} = 27$. \textit{(Làm tròn đến chữ số hàng phần trăm)}\\

\noindent \textbf{Câu 4:} Thuỷ phân hoàn toàn $14{,}6$ gam dipeptide Gly-Ala bằng dung dịch $\text{NaOH}$ dư, đun nóng. Cô cạn dung dịch sau phản ứng, khối lượng muối của các amino acid thu được là bao nhiêu gam? Cho nguyên tử khối: $\text{H} = 1$; $\text{C} = 12$; $\text{N} = 14$; $\text{O} = 16$; $\text{Na} = 23$. \textit{(Làm tròn đến chữ số hàng phần mười)}\\

\noindent \textbf{Câu 5:} Một loại phân đạm urea chứa $95\%$ khối lượng $(\text{NH}_2)_2\text{CO}$, phần còn lại là tạp chất không chứa nitrogen. Để cung cấp $70$ kg nguyên tố nitrogen cho một hecta lúa, người nông dân cần bón bao nhiêu kilôgam loại phân đạm này? Cho nguyên tử khối: $\text{H} = 1$; $\text{C} = 12$; $\text{N} = 14$; $\text{O} = 16$. \textit{(Làm tròn đến chữ số hàng đơn vị)}\\

\noindent \textbf{Câu 6:} Cho phương trình nhiệt hoá học của phản ứng đốt cháy methane:\\
$\text{CH}_4(g) + 2\text{O}_2(g) \longrightarrow \text{CO}_2(g) + 2\text{H}_2\text{O}(l)$\\
Biết enthalpy tạo thành chuẩn $\Delta_f H^{\circ}_{298}$ (kJ/mol) của $\text{CH}_4(g)$, $\text{CO}_2(g)$ và $\text{H}_2\text{O}(l)$ lần lượt là $-74{,}9$; $-393{,}5$ và $-285{,}8$. Đốt cháy hoàn toàn $4{,}0$ gam methane ở điều kiện chuẩn thì lượng nhiệt toả ra là bao nhiêu kJ? Cho nguyên tử khối: $\text{H} = 1$; $\text{C} = 12$. \textit{(Làm tròn đến chữ số hàng đơn vị)}\\\vspace{0.6cm}
\begin{center}
{\bfseries -{}-{}-{}-{}-{}-}\ \textbf{HẾT}\ {\bfseries -{}-{}-{}-{}-{}-}
\end{center}
\vspace{0.15cm}
{\small\itshape
Thí sinh không được sử dụng tài liệu. Cán bộ coi thi không giải thích gì thêm.}

\end{document}
